% TODO: re-structure hierarchy. right now everything is flat.

\section{Our Complete Dataset Collection}\label{app:rest}

We describe the rest of the dataset collection.

% Things to write.
% How it was created.
% Some statistics. See other dataset papers.
% Use cases
% Training potential
% Quality assurance / assessment

\subsection{\textsc{SightationVQA}}\label{app:gqa}
% \section{Generated Questions and Answers pairs}

In constructing $\textbf{Desc}_{\texttt{++}}$ for comparison with $\textbf{Desc}_{}$, we discovered that the quality of the Question–Answer pairs directly determines the quality of the resulting context. To clarify why we invested significant effort in carefully designing these question answer pairs, we employed an LLM as a judge to evaluate and classify them according to different quality levels. To measure the quality of the Question Answer pairs, we used the VLM-as-a-Judge prompt using GPT-4o model. The prompt itself is found in Appendix~\ref{app:prompts}.


% \begin{figure*}[htbp!]
%   \
%   \includegraphics[width=0.5\linewidth]{figures/evaluate_generated_qa.pdf} \hfill
%   \includegraphics[width=0.5\linewidth]{figures/evaluate_origin_qa.pdf}
%   \caption {Percentage Distribution of AI2D \textbf{generated (ours, left)} and \textbf{origin (right)} QA Scores}
%   \label{fig:evaluate_generated_qa}
% \end{figure*}

\begin{figure*}[h!]
  \centering
  \includegraphics[width=0.45\linewidth]{figures/interleaved_evaluation.png}
  \caption {Percentage distribution of the quality of question-answer pairs in AI2D and \textsc{SightationVQA}}
  \label{fig:evaluate_generated_qa}
\end{figure*}


Following \citet{chen2023alpagasus} and \citet{xu2024magpie}, we compared two sets of QA pairs with \textsc{GPT-4o}. Our generated QA sets are with up to six QA pairs for each of 4,903 diagrams, producing a total of 29,438 QA pairs (sometimes exceeding six pairs per diagram). As can be seen Figure \ref{fig:evaluate_generated_qa},  we found that 92.66\% of these our generated QA pairs were rated ``Excellent'', while 4.47\% were deemed ``Good'', underscoring their high quality. By contrast, the QA pairs sourced from the AI2D dataset, though numerous, included a large portion of masked or minimally informative queries. After filtering out these masked questions, we were left with 9,708 self-contained questions spanning 3,099 diagrams, where 73.86\% received an ``Excellent'' rating and 13.65\% were deemed ``Good''. This comparison reveals that our generated QA pairs provide a more robust and contextually relevant foundation, reinforcing the value of our meticulous QA design in constructing effective $\textbf{Desc}_{\texttt{++}}$.


\subsection{\textsc{SightationReasoning}}
Employing $\textbf{Desc}$ and $\textbf{Desc}_{\texttt{++}}$, we constructed \textsc{SightationReasoning}, a reasoning dataset that consists of reasoning path and reasoning QA pairs. The prompts used for the construction of reasoning datasets are found in Appendix~\ref{app:prompts}. To verify the quality of contents as a reasoning dataset, 10\% of the samples were randomly selected to be manually inspected.

\paragraph{Reasoning Path} The reasoning path explains the logical flow or deployment of the contents in a diagram such as cause-effect relationships, step-by-step processes, explanations of phenomena, comparions of contrasts, or dependencies between components. Employing 1k diagram images and descriptions in \textsc{Sightation}, the reasoning path was identified and generated by QVQ-72B-Preview. The reasoning path extracted from $\textbf{Desc}$ and $\textbf{Desc}_{\texttt{++}}$ is denoted as $\textbf{RPath}$ and $\textbf{RPath}_{\texttt{++}}$ respectively. Consequently, one diagram possesses two reasoning paths, resulting in 2k paths in total. 

\paragraph{Reasoning QA}
The reasoning QA encompasses five types of QA pairs that require a logical understanding of diagram contents and reasoning capabilities: Causal, Process, Conditional, Explanatory, and Reverse. Similarly to the reasoning path data, $\textbf{RQA}$ and $\textbf{RQA}_{\texttt{++}}$ were generated by \textsc{QVQ-72B-Preview} using 1k diagram images and descriptions. As a result, one diagram contains 10 reasoning QA pairs in which $\textbf{RQA}$ and $\textbf{RQA}_{\texttt{++}}$ respectively include 5 pairs. While \textsc{SightationVQA} covers the visual structure and details of a diagram, the reasoning QA in \textsc{SightationReasoning} consists of more knowledge-intensive questions that require logical thinking, paving the way for the reasoning applications of \textsc{Sightation}. 

\paragraph{Evaluation} The reasoning path of \textsc{SightationReasoning} can be used as an overall representation of "logical flow" or "relationships between instances" in a diagram when understanding it, which was emphasized in the BLV educator questionnaire. To make a model employ this information when responding to reasoning questions and evaluate the reasoning paths, we fed \textsc{Qwen2-VL-7B-Instruct} with $\textbf{RPath}$ and $\textbf{RPath}_{\texttt{++}}$ separately and asked it to solve 10 questions in $\textbf{RQA}$ and $\textbf{RQA}_{\texttt{++}}$. The similarity score between the gold answers and generated answers was calculated using BERTSCore \cite{zhang2019bertscore}, and the scores for the two cases both resulted in 0.975, verifying the equal usefulness of $\textbf{RPath}$ and $\textbf{RPath}_{\texttt{++}}$.


\section{Further Related Work}

In Table~\ref{tab:dataset_comparison_extended}, we extend Table~\ref{tab:dataset_comparison} for a more comprehensive view of neighboring datasets. To the best of our knowledge, there exists no dataset to date surpassing our contribution in terms of the breadth of use cases and granularity of validation with BLV individuals.

\begin{table*}[t]
    \centering
    \small
    \definecolor{blue}{RGB}{16,85,154}
    \definecolor{pink}{RGB}{219, 76, 119}
    \newcommand{\checkblue}{\textcolor{blue}{\checkmark}}
    \newcommand{\checkpink}{\textcolor{pink}{\texttimes}}
    
    \begin{tabularx}{\textwidth}{p{4.3cm} p{1.7cm} p{1.6cm} p{2.2cm} X}
        \toprule
        \makecell[c]{\textbf{Dataset}} & \textbf{\makecell[c]{Average \\Text \\Length}} & \textbf{\makecell[c]{Validated \\by\\ BLV?}} & \makecell[c]{\textbf{Applications}} & \textbf{\makecell[c]{Dimensions\\Assessed}}\\
        \midrule
        \makecell[l]{\textbf{\textsc{Sightation} (Ours)}\\
            \hspace{0.6cm}\textsc{-Completions}\\
            \hspace{0.6cm}\textsc{-Preference}\\
            \hspace{0.6cm}\textsc{-Retrieval}\\
            \hspace{0.6cm}\textsc{-VQA}\\
            \hspace{0.6cm}\textsc{-Reasoning}
        } & \makecell[r]{\textbf{188.3}\\(words)} & \makecell[c]{\textbf{\checkblue}} & \makecell[l]{$\cdot$ \textbf{Completion} \\ $\cdot$ \textbf{Preference}\\~~~\textbf{alignment} \\$\cdot$ \textbf{Retrieval} \\ $\cdot$ \textbf{Reward} \\~~~\textbf{modeling}} & \makecell[l]{$\cdot$ \textbf{Factuality} \\ $\cdot$ \textbf{Informativeness} \\ $\cdot$ \textbf{Succinctness} \\ $\cdot$ \textbf{Diversity} \\ $\cdot$ \textbf{Usefulness}, \\~~~in 4 finer aspects \\ $\cdot$ \textbf{Interpretiveness}} \\
        \midrule
        VisText \cite{tang2023vistext} & \makecell[r]{74.6} & \makecell[c]{\checkpink} & Completion & Accuracy, Descriptiveness \\
        MathVista \cite{lu2023mathvista} & \makecell[r]{58.0} & \makecell[c]{\checkpink} & VQA, Reasoning & Correctness \\
        ChartGemma \cite{masry2024chartgemma} & \makecell[r]{37.5} & \makecell[c]{\checkpink} & Completion & Informativeness, Factual Cor-
rectness, Structure \\
        CBD \cite{bhushan-lee-2022-block} & \makecell[r]{114.5} & \makecell[c]{\checkpink} & Summarization & Adequacy, Fluency, Coherence \\
        VizWiz-VQA \cite{gurari2018vizwiz} & \makecell[r]{8.6} &  \makecell[c]{\checkblue} & VQA & Diversity, Answerability \\
        VizWiz-LF \cite{huh2024long} & \makecell[r]{73.2} & \makecell[c]{\checkblue} & VQA & Relevance, Helpfulness, Plausi-
bility, Fluency, Correctness \\
        DiagramQG \cite{zhang2024diagramqg} & \makecell[r]{9.5} & \makecell[c]{\checkpink} & DQA & Diversity, Object Density \\
        ScienceQA \cite{lu2022learn} & \makecell[r]{119.7} & \makecell[c]{\checkpink} & VQA, Reasoning & Correctness \\
        ChartQA \cite{masry-etal-2022-chartqa} & \makecell[r]{13.0} & \makecell[c]{\checkpink} & VQA & Syntactic Diversity \\
        Flickr8K \cite{hodosh2013framing} & \makecell[r]{11.8} & \makecell[c]{\checkpink} & Description & Diversity \\
        PASCAL-50S \cite{vedantam2015cider} & \makecell[r]{8.8} & \makecell[c]{\checkpink} & Description & Factuality, Literality, Generality \\
        Polaris \cite{wada2024polos} & \makecell[r]{11.5} & \makecell[c]\checkpink & Description & Fluency, Relevance, Descriptiveness \\
        %Visual-CoT \cite{shao2024visual} & \makecell[r]{TBD} & \makecell[c]{\checkpink} & VQA, Reasoning & TBD \\
        %VQAOnline \cite{chen2024fully} & \makecell[r]{TBD} & \makecell[c]{\checkpink} & VQA & TBD \\
        Multimodal Arxiv \cite{li2024multimodal} & \makecell[r]{49.7} & \makecell[c]{\checkpink} & Description, VQA, Reasoning & Factual Alignment, Visual Clarity, Unambiguous Textual Information, Question and Option Relevance, Comprehensive Integration, Equitable Content \\
        MMMU \cite{Yue_2024_CVPR} & \makecell[r]{53.2} & \makecell[c]{\checkpink} & VQA, Reasoning & Difficulty, Knowledge, Reasoning \\
        \bottomrule
    \end{tabularx}
    \caption{Extended related work.}
    \label{tab:dataset_comparison_extended}
\end{table*}

\begin{figure}[t]
    \centering
    \includegraphics[width=0.45\linewidth]{figures/abstract.pdf}
    \caption{Less can be more for BLV users. Our approach streamlines details to highlight the core information while emphasizing key details to increase information density and maximize information efficiency per unit length.}
    \label{fig:abstract_img}
\end{figure}



\begin{table*}[htbp!]
\centering

\begin{tabular}{ll}
\toprule
Notation & Description \\ \cmidrule{1-2}
$(\cdot)^{model}$      & \pbox{0.7\textwidth}{The description \textbf{Desc} generated by (or an annotation on a generation from) a $model \in \{\texttt{g}, \texttt{q}\}$, for \textsc{GPT-4o mini} and \textsc{Qwen2-VL}, respectively. Later overloaded with narrower descriptors, such as \texttt{base}, \texttt{sft}, and \texttt{sft+dpo} to refer to the baseline/tuned models.} \\ \cmidrule{1-2}

$(\cdot)_{anchor}$      & \pbox{0.7\textwidth}{The conditioning input at the description generation stage. $anchor \in \{\texttt{None}, \texttt{++}\}$, for the one-pass image-only conditioning and the two-pass image+QA conditioning, respectively.}  \\ \cmidrule{1-2}

$\textbf{Preference}^{model}$  & \pbox{0.7\textwidth}{Preference annotation between two $\textbf{Desc}^{model}$'s on different conditioning inputs. Value takes either of the $anchor$ set \{\texttt{None}, \texttt{++}\}}  \\ \cmidrule{1-2}

$Aspect^{model}_{anchor}$      & \pbox{0.7\textwidth}{Rating annotation in terms of $Aspect \in$ \{\textbf{Factuality}, \textbf{Informativeness}, \textbf{Succinctness}, \textbf{Diversity}, \textbf{Usefulness-Gen}, \textbf{Usefulness-Sum}, \textbf{Usefulness-MCQ}, \textbf{Usefulness-OEQ}, \textbf{Nature}\}, for a description generated by $model$ conditioned on $anchor$. Value is an integer ranging from 1 to 5, on the 5-point Likert scale.}   \\ \cmidrule{1-2}

$\textbf{Best}^{model}_{anchor}$ & \pbox{0.7\textwidth}{Best sentence annotation. Value is a substring of $\textbf{Desc}^{model}_{anchor}$.}  \\

\bottomrule
\end{tabular}
\caption{Notations}
\label{tab:notation}
\end{table*}

\section{Details on the Annotations}\label{app:ann_detail}

\subsection{Logistics}

All experimentation was reviewed and approved by the Institutional Review Board. Recruiting the sighted general group was done via an online forum. Each sighted general group annotator was paid an approximate equivalent of USD80 for completing the assigned task. Recruiting the educators was done by directly corresponding with the schools for the blind. A sighted educator was compensated an approximate equivalent of USD80. A BLV educator was compensated an approximate equivalent of USD80 to USD160, depending on the number of samples completed.

\subsection{Annotations Statistics}
\paragraph{Preliminaries}

Of the 1,000 diagrams distributed to the annotators, 956 have been annotated by three annotators; 41 by two; 1 by a single annotator; and 2 by none. We collected annotations on 3,992 diagram-description pairs, each with at most 3 annotations.

% TODO: include histogram, all sorts

% TODO: include wins from Sighted Educators and BLV Educators (perhaps in the body)
\paragraph{Internal Consistency}

In Table~\ref{tab:cronbach}, we report the Cronbach's alpha value for each assessment group. The statistic is widely interpreted as the reliability of a set of survey items.

\begin{table*}[h!]
\centering

\begin{tabular}{lr}
\toprule
\textbf{Group}                   & \textbf{Cronbach's $\alpha$} \\ \cmidrule{1-2}
Sighted General         & 0.70 \\
Sighted Educators       & 0.94 \\
BLV Educators           & 0.80 \\

\bottomrule
\end{tabular}
\caption{Our survey items are considered of acceptable ($\ge0.7$) to excellent ($\ge0.9$) reliability.}
\label{tab:cronbach}
\end{table*}


\paragraph{Point-Biserial Correlation}

We examine the relationship between the binary variable, \textbf{Preference}, and the 5-point scale ratings per aspect.

\begin{table*}[h!]
\centering

\begin{tabular}{lrrrrr}
\toprule
                        & \multicolumn{5}{c}{\textbf{Aspects}} \\ \cmidrule{2-6}
\textbf{Group}          & \textbf{Factuality} & \textbf{Informativeness} & \textbf{Succinctness} & \textbf{Diversity} & \textbf{Usefulness-Gen} \\ \cmidrule{1-6}
Sighted General         & $0.36^{***}$ & $0.37^{***}$ & $0.31^{***}$ & $0.34^{***}$ & $0.43^{***}$ \\
Sighted Educators       & $0.25^{***}$ & $0.30^{***}$ & $0.30^{***}$ & $0.34^{***}$ & --- \\
\bottomrule
\end{tabular}
\caption{Correlation values between preference choice and aspect ratings were found to be moderately positive and statistically significant. (***: $p < 0.001$)}
\label{tab:point_biserial}
\end{table*}


\paragraph{Cohen's $d$}

Cohen's $d$ is a widely used statistic to measure the size of the effect of a treatment. It is the difference in the means of the treatment and control groups, normalized by the pooled standard deviation. By guidelines set forth by Cohen himself, values over 0.2 are typically considered a small effect size; 0.5, medium; and 0.8, large.





\begin{figure*}[htbp!]
    \centering
    \includegraphics[width=0.45\textwidth]{figures/wins_average.png}
    \caption{Win rates by $model$.}
    \label{fig:overall_wins}
\end{figure*}


\subsection{Annotations Post-processing}

\paragraph{Preference Choice}
We aggregate the multiple annotations on the basis of majority. That is, for the three-annotation samples, a 3:0 or 2:1 is considered a ``victory'' and the victor \textbf{Desc} wins that sample. For two-annotation samples with differing preferences, a tie is recorded. The overall win-loss statistics normalized against the number of diagrams (998) is shown in Figure~\ref{fig:overall_wins}.


\paragraph{Rating Assessment}

% A summary of sighted users' assessment is presented in Figures~\ref{fig:overall_fac}-\ref{fig:overall_div}. Each figure shows the preference choice counts, regardless of the score difference. For instance, a rating of $\textbf{Factuality}^{\texttt{g}}_{} = 5$ versus a rating of $\textbf{Factuality}^{\texttt{g}}_{\texttt{++}} = 1$ is considered 1 win for $\textbf{Desc}^{\texttt{g}}_{}$, just as the case where the ratings are 3 vs. 2.


\subparagraph{Best Sentence Choice}

The best sentence for each context was manually selected by BLV annotators after listening to the context. We analyzed people’s preferences by examining the position and length of the best sentence within each context.

\subparagraph{Position}

The normalized position of the best sentence is shown in Figures~\ref{fig:best_sentence_distribution_g}-\ref{fig:best_sentence_distribution_q}. To calculate the relative position, both the context and the best sentence were tokenized at the word level, and the position of the overlapping best sentence within the context was identified. This position was then normalized to a value between 0 and 1 by dividing it by the total length of the context. Furthermore, since some BLV annotators could not select a best sentence within the context, a filtering step was applied by setting an overlap threshold of 0.9 to account for such cases.

Figures~\ref{fig:best_sentence_distribution_g}-\ref{fig:best_sentence_distribution_q} illustrate that the best sentences in each context are predominantly positioned at the beginning and end. This pattern can be attributed to cognitive biases, specifically primacy bias and recency bias. Primacy bias refers to the tendency to place greater importance on the first pieces of information encountered in a sequence, while recency bias reflects the tendency to prioritize the most recently encountered information. Consequently, these biases increase the likelihood that preferred sentences will be selected from the beginning and end of the context.


\begin{figure*}[htbp!]
  \
  \includegraphics[width=0.48\linewidth]{figures/best_sentence_distribution_ag.png} \hfill
  \includegraphics[width=0.48\linewidth]{figures/best_sentence_distribution_bg.png}
  \caption {Descriptions generated by \textsc{GPT-4o mini}}
  \label{fig:best_sentence_distribution_g}
\end{figure*}

\begin{figure*}[htbp!]
  \
  \includegraphics[width=0.48\linewidth]{figures/best_sentence_distribution_aq.png} \hfill
  \includegraphics[width=0.48\linewidth]{figures/best_sentence_distribution_bq.png}
  \caption {Descriptions generated by \textsc{Qwen2-VL}}
  \label{fig:best_sentence_distribution_q}
\end{figure*}



\subparagraph{Length}

The length of the best sentence in each context is presented in Figure~\ref{fig:best_sentence_length_distribution}. The length was determined by counting the total number of words in the best sentence. As shown in Figure 10, the best sentences across different contexts predominantly consist of 20 to 30 words, exhibiting a similar distribution pattern.

\begin{figure}[htbp!]
    \centering
    \includegraphics[width=0.48\linewidth]{figures/best_sentence_length_word.png}
    \caption{boxplot of best sentence length}
    \label{fig:best_sentence_length_distribution}
\end{figure}


\section{Retrieval Dataset Construction}\label{app:retrieval}

The winner among the four human-annotated descriptions was assigned as the top 1 positive in terms of preference and average rating. The top 5 set contains all 4 human-annotated descriptions and 1 synthesized description; the top 10 set is a superset of the top 5, joined by 5 more synthetic descriptions. The synthetic descriptions are perturbed versions of the human-annotated descriptions, each missing a random, non-best sentence. The 10 hard negatives for an image were selected among the combined pool of top 1 descriptions for other images, sorted by cosine similarity in the embedding space. The embeddings were computed by a widely used sentence transformer, \textsc{all-mpnet-base-v2}\citep{song2020mpnet}.



\section{Detailed Results}\label{app:results}

We report the VLM-as-a-Judge evaluation and classic metric results in Tables~\ref{tab:big1a}, ~\ref{tab:big1b}, ~\ref{tab:big2B}, and ~\ref{tab:big7B}.

\subsection{Evaluation by Automatic Metrics}

% TODO: include correlation with all three groups.
% argue which group it most strongly behaves like.
% TODO: describe our VLM judge in the section where classic metrics are mentioned.
% -> Done.
\paragraph{\textsc{QVQ-72B-Preview}}

On GPT and Qwen 72B generations, the VLM judge did not reveal significant difference between the two anchors, and the little differences present aligned with assessments by the sighted general group, as can be expected from a general-purpose VLM.

It is important to note that even a state-of-the-art VLM fails to capture the BLV perspectives in text evaluation. 
% TODO: that's why we did some few-shot prompting with them and the correlation is...

% TODO: refer to the list of metrics mentioned earlier (label that section and then ref)
\paragraph{Classic Metrics}

To our surprise, almost all instances of classic metric evaluations resulted in a win for the \texttt{++} anchor. However, the numbers from classic metrics evaluation are more of a shortcoming on the part of the classic metrics, rather than an accurate portrayal of the effectiveness of our proposed latent supervision. This is because our ``gold'' ground truths from BLV educators show that, while the QA-guided generation does manifest in ways beneficial to BLV individuals, classic automatic metrics poorly represent the assessment space covered by BLV, such as with the \textbf{Diversity} and \textbf{Usefulness-OEQ} aspects.



\begin{table*}
    \centering
    \begin{tabular}{llll}
    \toprule
    Experiment ID & & \multicolumn{2}{c}{\textbf{Assessments for}} \\ \cmidrule{3-4}
    \textbf{Description Generators}         & \textbf{Metrics} & $\textbf{Desc}_{}$ & $\textbf{Desc}_{\texttt{++}}$ \\ \midrule
    \multirow{28}{0.2\textwidth}{Experiment 1a\\
    \textsc{GPT-4o mini}\\vs.\\\textsc{GPT-4o mini}}
    
        & CLIP Score                                    & 0.476 & \textbf{0.524} \\ 
        & SigLIP Score                                  & \textbf{0.921} & 0.914 \\ 
        & BLIP-2 Retrieval Score                        & 0.495 & \textbf{0.505} \\ 
        & Self-BLEU                                     & 0.256 & \textbf{0.268} \\ 
        & PAC-Score                                     & 0.699 & \textbf{0.703} \\ 
        & LongCLIP-B Score                                & \textbf{0.507} & 0.493 \\
        & LongCLIP-L Score                                & \textbf{0.531} & 0.469 \\\cmidrule{2-4}                                    

        & $\cdot$ VLM-as-a-Judge Evaluation Average     & \textbf{4.080} & 4.033 \\
        & Factuality                                    & 4.433 & \textbf{4.445} \\
        & Informativeness                               & \textbf{4.200} & 4.166 \\
        & Succinctness                                  & 4.108 & \textbf{4.146} \\
        & Diversity                                     & \textbf{3.578} & 3.375 \\\cmidrule{2-4}

        & $\cdot$ Sighted General Group Average         & \textbf{3.983} & 3.962 \\
        & Factuality                                    & \textbf{4.128} & 4.093 \\
        & Informativeness                               & \textbf{4.367} & 4.032 \\
        & Succinctness                                  & 3.556 & \textbf{4.040} \\
        & Diversity                                     & \textbf{3.879} & 3.685 \\

        % might get rid of the two groups below -> depending on whether annotators respond.
        & $\cdot$ Sighted Educator Group Average        & 3.22 & \textbf{3.35} \\
        & Factuality                                    & \textbf{3.35} & 3.30 \\
        & Informativeness                               & 3.43 & 3.43 \\
        & Succinctness                                  & 2.78 & \textbf{3.53} \\
        & Diversity                                     & \textbf{3.18} & 3.08 \\
        & Usefulness to BLV                             & 3.35 & \textbf{3.40} \\
        
        & $\cdot$ BLV Educator Group Average            & 2.98 & \textbf{3.17} \\
        & Succinctness                                  & 2.43 & \textbf{2.55} \\
        & Diversity                                     & \textbf{3.23} & 3.15 \\
        & Usefulness, Summary                           & 2.95 & \textbf{3.33} \\
        & Usefulness, Multiple-chioce Questions         & 3.20 & \textbf{3.28} \\
        & Usefulness, Open-ended Questions              & 2.88 & \textbf{3.13} \\
        & Nature of Context                             & 2.98 & 3.17 \\

    \bottomrule
    \end{tabular}
\caption{The full evaluation on descriptions by GPT. Nature of Context values are not in bold because it is a categorical variable.}\label{tab:big1a}
\end{table*}

\begin{table*}
    \centering
    \begin{tabular}{llll}
    \toprule
    Experiment ID & & \multicolumn{2}{c}{\textbf{Assessments for}} \\ \cmidrule{3-4}
    \textbf{Description Generators}         & \textbf{Metrics} & $\textbf{Desc}_{}$ & $\textbf{Desc}_{\texttt{++}}$ \\ \midrule
        
    \multirow{28}{0.2\textwidth}{Experiment 1b\\
    \textsc{Qwen2-VL-72B-Instruct}\\vs.\\\textsc{Qwen2-VL-72B-Instruct}}
    
        & CLIP Score                                    & 0.451 & \textbf{0.549} \\ 
        & SigLIP Score                                  & 0.911 & \textbf{0.932} \\ 
        & BLIP-2 Retrieval Score                        & 0.494 & \textbf{0.506} \\ 
        & Self-BLEU                                     & 0.260 & \textbf{0.274} \\ 
        & PAC-Score                                     & 0.709 & \textbf{0.716} \\
        & LongCLIP-B                                    & 0.443 & \textbf{0.610} \\
        & LongCLIP-L                                    & 0.468 & \textbf{0.532} \\\cmidrule{2-4}

        & $\cdot$ VLM-as-a-Judge Evaluation Average     & \textbf{4.094} & 3.916 \\
        & Factuality                                    & \textbf{4.483} & 4.428 \\
        & Informativeness                               & \textbf{4.239} & 3.952 \\
        & Succinctness                                  & 4.026 & \textbf{4.072} \\
        & Diversity                                     & \textbf{3.629} & 3.210 \\ \cmidrule{2-4}

        & $\cdot$ Sighted General Group Average         & \textbf{4.002} & 3.850 \\
        & Factuality                                    & 3.982 & \textbf{4.060} \\
        & Informativeness                               & \textbf{4.233} & 3.782 \\
        & Succinctness                                  & 3.889 & \textbf{4.035} \\
        & Diversity                                     & \textbf{3.905} & 3.523 \\

        & $\cdot$ Sighted Educator Group Average        & 4.01 & \textbf{4.13} \\
        & Factuality                                    & 4.05 & 4.05 \\
        & Informativeness                               & \textbf{4.38} & 4.13 \\
        & Succinctness                                  & 3.80 & \textbf{4.48} \\
        & Diversity                                     & 3.80 & \textbf{3.83} \\
        & Usefulness to BLV                             & 4.03 & \textbf{4.15} \\

        % & $\cdot$ BLV Educator Group Average            & & \\
        % & Succinctness                                  & & \\
        % & Diversity                                     & & \\
        % & Usefulness, Summary                           & & \\
        % & Usefulness, Multiple-chioce Questions         & & \\
        % & Usefulness, Open-ended Questions              & & \\
        % & Nature of Context                             & & \\
     

    \bottomrule
    \end{tabular}
\caption{The full evaluation on descriptions by the 72B model. Due to limited recruiting, BLV annotators were not given this set.}\label{tab:big1b}
\end{table*}

% might exclude DPO results
% as a dataset paper, showing SFT+DPO solely for context generation seems less important than
% illustrating use cases in, say, reasoning, retrieval, ranking, and/or rag can be more appealing

\begin{landscape}
\begin{table}
    \small
    \centering
    \captionsetup{width=\textwidth}
    \begin{tabular}{l>{\columncolor{gray!15}}l>{\columncolor{gray!15}}lll>{\columncolor{gray!15}}l>{\columncolor{gray!15}}l}
    \toprule
    Fine-tuning \textcolor{ColorBlue}{\textsc{Qwen2-VL-2B-Instruct}}& \multicolumn{6}{c}{\textbf{Pairwise Assessments for}  $\textbf{Desc}^{\textcolor{ColorBlue}{\texttt{q2b}}}_{}$ vs. $\textbf{Desc}^{\textcolor{ColorBlue}{\texttt{q2b}}}_{\texttt{++}}$ } \\ \cmidrule{2-7}

    \textbf{Metrics (Scores) by} & $\textbf{Desc}^{\texttt{base}}_{}$ & $\textbf{Desc}^{\texttt{base}}_{\texttt{++}}$ & $\textbf{Desc}^{\texttt{sft}}_{}$ & $\textbf{Desc}^{\texttt{sft}}_{\texttt{++}}$ & $\textbf{Desc}^{\texttt{sft+dpo}}_{}$  & $\textbf{Desc}^{\texttt{sft+dpo}}_{\texttt{++}}$ \\ \midrule

        CLIP Score                              & 0.442 & \textbf{0.558} & 0.466 & \textbf{0.534} & 0.451 & \textbf{0.549} \\ 
        SigLIP Score                            & 0.916 & \textbf{0.941} & 0.911 & \textbf{0.931} & 0.914 & \textbf{0.940} \\ 
        BLIP-2 Retrieval Score                  & 0.491 & \textbf{0.509} & 0.493 & \textbf{0.507} & 0.491 & \textbf{0.509} \\ 
        Self-BLEU                               & 0.274 & \textbf{0.278} & 0.285 & \textbf{0.291} & 0.277 & \textbf{0.281} \\ 
        PAC-Score                               & 0.711 & \textbf{0.718} & 0.706 & \textbf{0.710} & 0.712 & \textbf{0.718} \\
        LongCLIP-B                              & 0.419 & \textbf{0.581} & 0.452 & \textbf{0.548} & 0.445 & \textbf{0.555} \\ 
        LongCLIP-L                              & 0.417 & \textbf{0.583} & 0.454 & \textbf{0.546} & 0.459 & \textbf{0.541} \\ \cmidrule{1-7}     

        $\cdot$ VLM-as-a-Judge Evaluation Average & 3.307 & \textbf{3.509} & \textbf{3.732} & 3.663 & 3.334 & \textbf{3.519} \\
        Factuality                              & 3.426 & \textbf{3.783} & 3.926 & \textbf{3.974} & 3.431 & \textbf{3.784} \\
        Informativeness                         & 3.394 & \textbf{3.567} & \textbf{3.854} & 3.715 & 3.438 & \textbf{3.577} \\
        Succinctness                            & 3.346 & \textbf{3.662} & 3.707 & \textbf{3.774} & 3.347 & \textbf{3.659} \\
        Diversity                               & \textbf{3.062} & 3.025 & \textbf{3.442} & 3.188 & \textbf{3.118} & 3.054 \\ \cmidrule{1-7}

        $\cdot$ Sighted Educators Group Average & 3.91 & \textbf{3.95} & & & 4.34 & \textbf{4.49} \\
        Factuality                              & 3.95 & \textbf{4.03} & & & 4.42 & \textbf{4.66} \\
        Informativeness                         & 4.03 & \textbf{4.05} & & & 4.39 & \textbf{4.50} \\
        Succinctness                            & \textbf{3.98} & 3.90 & & & 4.37 & \textbf{4.50} \\
        Diversity                               & 3.65 & \textbf{3.80} & & & 4.18 & \textbf{4.32} \\
        Usefulness to BLV                       & 3.93 & \textbf{3.98} & & & 4.34 & \textbf{4.50} \\

        $\cdot$ BLV Educators Group Average     & \textbf{3.33} & 3.25 & \multicolumn{2}{c}{---} & 2.62 & \textbf{3.17} \\
        Succinctness                            & \textbf{3.45} & 3.33 & & & 3.15 & \textbf{3.30} \\
        Diversity                               & \textbf{3.18} & 3.10 & & & 2.03 & \textbf{2.53} \\
        Usefulness, Summary                     & \textbf{3.53} & 3.40 & & & 2.88 & \textbf{3.45} \\
        Usefulness, Multiple-choice Questions   & \textbf{3.15} & 3.10 & & & 2.88 & \textbf{3.73} \\
        Usefulness, Open-ended Questions        & 3.15 & \textbf{3.21} & & & 2.28 & \textbf{3.00} \\
        Nature of Context                       & 3.33 & 3.25 & & & 2.50 & 3.00 \\
                                    
    \bottomrule
    \end{tabular}
\caption{Evaluation of the 2B model from baseline to SFT to DPO. Note that human evaluation results are unnormalized values on the 5-point Likert scale, so direct comparisons are meaningful only within the pairwise shaded columns. SFT versus SFT samples were not distributed due to limited annotator resources. Nature of Context values are not in bold because it is a categorical variable.}\label{tab:big2B}
\end{table}
\end{landscape}

\begin{landscape}
\begin{table}
    \small
    \centering
    \captionsetup{width=\textwidth}
    \begin{tabular}{l>{\columncolor{gray!15}}l>{\columncolor{gray!15}}lll>{\columncolor{gray!15}}l>{\columncolor{gray!15}}l}
    \toprule
    Fine-tuning \textcolor{ColorOrange}{\textsc{Qwen2-VL-7B-Instruct}}& \multicolumn{6}{c}{\textbf{Pairwise Assessments for}  $\textbf{Desc}^{\textcolor{ColorOrange}{\texttt{q7b}}}_{}$ vs. $\textbf{Desc}^{\textcolor{ColorOrange}{\texttt{q7b}}}_{\texttt{++}}$ } \\ \cmidrule{2-7}

    \textbf{Metrics (Scores) by} & $\textbf{Desc}^{\texttt{base}}_{}$ & $\textbf{Desc}^{\texttt{base}}_{\texttt{++}}$ & $\textbf{Desc}^{\texttt{sft}}_{}$ & $\textbf{Desc}^{\texttt{sft}}_{\texttt{++}}$ & $\textbf{Desc}^{\texttt{sft+dpo}}_{}$  & $\textbf{Desc}^{\texttt{sft+dpo}}_{\texttt{++}}$ \\ \midrule

        CLIP Score                              & 0.423 & \textbf{0.577} & 0.411 & \textbf{0.589} & 0.407 & \textbf{0.593} \\ 
        SigLIP Score                            & 0.922 & \textbf{0.952} & 0.918 & \textbf{0.944} & 0.923 & \textbf{0.952} \\ 
        BLIP-2 Retrieval Score                  & 0.490 & \textbf{0.510} & 0.489 & \textbf{0.511} & 0.490 & \textbf{0.510} \\ 
        Self-BLEU                               & 0.268 & \textbf{0.274} & 0.275 & \textbf{0.282} & 0.268 & \textbf{0.275} \\ 
        PAC-Score                               & 0.713 & \textbf{0.720} & 0.706 & \textbf{0.714} & 0.711 & \textbf{0.718} \\
        LongCLIP-B                              & 0.419 & \textbf{0.581} & 0.452 & \textbf{0.589} & 0.417 & \textbf{0.583} \\
        LongCLIP-L                              & 0.417 & \textbf{0.583} & 0.486 & \textbf{0.514} & 0.412 & \textbf{0.588} \\ \cmidrule{1-7}  

        $\cdot$ VLM-as-a-Judge Evaluation Average & \textbf{3.951} & 3.652 & \textbf{4.021} & 3.758 & \textbf{3.948} & 3.642\\
        Factuality                              & \textbf{4.271} & 4.157 & \textbf{4.371} & 4.261 & \textbf{4.289} & 4.161 \\
        Informativeness                         & \textbf{4.101} & 3.645 & \textbf{4.161} & 3.770 & \textbf{4.100} & 3.642 \\
        Succinctness                            & \textbf{3.946} & 3.892 & \textbf{3.974} & 3.964 & \textbf{3.904} & 3.858 \\
        Diversity                               & \textbf{3.486} & 2.913 & \textbf{3.576} & 3.036 & \textbf{3.498} & 2.906 \\ \cmidrule{1-7}

        % 샘플 갯수 40개 뿐이라 둘째자리에서 끊음. 어차피 0이나 5일 셋째자리 빼고.
        $\cdot$ Sighted Educators Group Average & \textbf{4.37} & 3.97 & & & \textbf{3.97} & 3.95 \\
        Factuality                              & \textbf{4.82} & 4.56 & & & \textbf{4.00} & 3.95 \\
        Informativeness                         & \textbf{4.67} & 3.87 & & & 4.08 & \textbf{4.13} \\
        Succinctness                            & 3.95 & \textbf{4.15} & & & 3.88 & \textbf{4.00} \\
        Diversity                               & \textbf{4.23} & 3.64 & & & \textbf{3.88} & 3.70 \\
        Usefulness to BLV                       & \textbf{4.37} & 3.97 & & & \textbf{4.03} & 3.95 \\

        $\cdot$ BLV Educators Group Average     & \textbf{3.87} & 3.82 & \multicolumn{2}{c}{---} & \textbf{3.82} & 3.71 \\
        Succinctness                            & 4.30 & \textbf{4.55} & & & 4.48 & \textbf{4.65} \\
        Diversity                               & 4.20 & 4.20 & & & \textbf{4.13} & 3.90 \\
        Usefulness, Summary                     & 4.15 & \textbf{4.55} & & & 4.25 & \textbf{4.35} \\
        Usefulness, Multiple-choice Questions   & \textbf{4.40} & 4.20 & & & \textbf{4.15} & 3.95 \\
        Usefulness, Open-ended Questions        & 3.80 & 3.80 & & & \textbf{3.70} & 3.58 \\
        Nature of Context                       & 2.35 & 1.60 & & & 2.23 & 1.85 \\
                                    
    \bottomrule
    \end{tabular}
\caption{Evaluation of the 7B model. Note that human evaluation results are nominal values on the 5-point Likert scale, so direct comparisons are meaningful only within the pairwise shaded columns. As with the 2B case, SFT versus SFT samples were not distributed due to limited annotator resources. Nature of Context values are not in bold because it is a categorical variable.}\label{tab:big7B}
\end{table}
\end{landscape}




% possibly include emulated blind VLM judges -> depending on BLV annotators' results


% TODO: Consider separating the human-annotated results. Because num_samples is different.
% classic and LLMJudge: 998
% sighted and BLV educators: 40
% depending on BLV annotations, might be removed to appendix or completely
\begin{table*}\label{tab:big3a}
    \centering
    \begin{tabular}{llll}
    \toprule
    Experiment ID & & \multicolumn{2}{c}{\textbf{Assessments for}} \\ \cmidrule{3-4}
    \textbf{Description Generators}         & \textbf{Metrics} & $\textbf{Desc}^{\textcolor{ColorBlue}{\texttt{q72bbase}}}_{}$ & $\textbf{Desc}^{\textcolor{ColorOrange}{\texttt{q7bdpo}}}_{\texttt{++}}$  \\ \midrule
    \multirow{20}{0.2\textwidth}{Experiment 3a\\\textsc{\textcolor{ColorBlue}{Qwen2-VL-72B-Instruct}}\\vs.\\\textsc{\textcolor{ColorOrange}{Fine-tuned Qwen2-VL-7B-Instruct}}} 
        & CLIP Score                                & 0.390 & \textbf{0.610} \\ 
        & SigLIP Score                              & 0.911 & \textbf{0.952} \\ 
        & BLIP-2 Retrieval Score                    & 0.487 & \textbf{0.513} \\ 
        & Self-BLEU                                 & 0.260 & \textbf{0.275} \\ 
        & PAC-Score                                 & 0.709 & \textbf{0.719} \\
        & LongCLIP-B Score                          & 0.388 & \textbf{0.612} \\
        & LongCLIP-L Score                          & 0.445 & \textbf{0.555} \\\cmidrule{2-4}                                  

        & $\cdot$ VLM-as-a-Judge Evaluation Average & \textbf{4.095} & 3.650 \\
        & Factuality                                & \textbf{4.477} & 4.238 \\
        & Informativeness                           & \textbf{4.262} & 3.586 \\
        & Succinctness                              & \textbf{3.990} & 3.894 \\
        & Diversity                                 & \textbf{3.652} & 2.880 \\\cmidrule{2-4}

        & $\cdot$ Sighted Educators Group Average   & \textbf{3.21} & 3.01 \\
        & Factuality                                & \textbf{3.30} & 3.28 \\
        & Informativeness                           & \textbf{3.33} & 2.95 \\
        & Succinctness                              & 2.95 & \textbf{3.18} \\
        & Diversity                                 & \textbf{3.13} & 2.68 \\
        & Usefulness to  BLV                        & \textbf{3.35} & 2.98 \\

        & $\cdot$ BLV Educators Group Average       & 3.69 & \textbf{4.33} \\
        & Succinctness                              & 3.60 & \textbf{4.55} \\
        & Diversity                                 & 3.60 & \textbf{3.90} \\
        & Usefulness, Summary                       & 3.95 & \textbf{4.30} \\
        & Usefulness, Multiple-choice Questions     & 3.70 & \textbf{4.55} \\
        & Usefulness, Open-ended Questions          & 3.70 & \textbf{4.45} \\
        & Nature of Context                         & 3.60 & 4.25 \\

    \bottomrule
    \end{tabular}
\caption{The smaller model outperforms a larger variant across many metrics. It is also important to note that the VLM judgments align better with sighted educators than with BLV educators. Further analysis is found in Section~\ref{sec:november}. This tendency is especially strong with the pairwise comparison between 72B- and 7B-generated descriptions. Nature of Context values are not in bold because it is a categorical variable.}
\end{table*}


\begin{table*}\label{tab:big3b}
    \centering
    \begin{tabular}{llll}
    \toprule
    Experiment ID & & \multicolumn{2}{c}{\textbf{Assessments for}} \\ \cmidrule{3-4}
    \textbf{Description Generators}         & \textbf{Metrics} & $\textbf{Desc}^{\textcolor{ColorBlue}{\texttt{q7bbase}}}_{}$ & $\textbf{Desc}^{\textcolor{ColorOrange}{\texttt{q2bdpo}}}_{\texttt{++}}$  \\ \midrule

    \multirow{20}{0.2\textwidth}{Experiment 3b \\\textsc{\textcolor{ColorBlue}{Qwen2-VL-7B-Instruct}}\\vs.\\\textsc{\textcolor{ColorOrange}{Fine-tuned Qwen2-VL-2B-Instruct}}}
    
        & CLIP Score                                & 0.486 & \textbf{0.514} \\ 
        & SigLIP Score                              & 0.922 & \textbf{0.940} \\ 
        & BLIP-2 Retrieval Score                    & 0.500 & 0.500 \\ 
        & Self-BLEU                                 & 0.268 & \textbf{0.281} \\ 
        & PAC-Score                                 & 0.713 & \textbf{0.718} \\
        & LongCLIP-B Score                          & 0.316 & \textbf{0.684} \\
        & LongCLIP-L Score                          & \textbf{0.559} & 0.441 \\\cmidrule{2-4}                                       

        & $\cdot$ VLM-as-a-Judge Evaluation Average & \textbf{3.921} & 3.545 \\
        & Factuality                                & \textbf{4.203} & 3.935 \\
        & Informativeness                           & \textbf{4.046} & 3.592 \\
        & Succinctness                              & \textbf{3.942} & 3.709 \\
        & Diversity                                 & \textbf{3.493} & 2.945 \\ \cmidrule{2-4}

        & $\cdot$ Sighted Educators Group Average   & \textbf{4.75} & 4.44 \\
        & Factuality                                & \textbf{4.75} & 4.50 \\
        & Informativeness                           & \textbf{4.65} & 4.38 \\
        & Succinctness                              & \textbf{4.88} & 4.40 \\
        & Diversity                                 & \textbf{4.80} & 4.63 \\
        & Usefulness to BLV                         & \textbf{4.65} & 4.28 \\

        & $\cdot$ BLV Educators Group Average       & 4.13 & \textbf{4.32} \\
        & Succinctness                              & 4.05 & \textbf{4.15} \\
        & Diversity                                 & 4.08 & \textbf{4.15} \\
        & Usefulness, Summary                       & 3.85 & \textbf{4.13} \\
        & Usefulness, Multiple-choice Questions     & 4.53 & \textbf{4.58} \\
        & Usefulness, Open-ended Questions          & 4.23 & \textbf{4.35} \\
        & Nature of Context                         & 4.08 & 4.50 \\

    \bottomrule
    \end{tabular}
\caption{The 2B model performs on par with the 7B variant. Again, VLM judgments align better with sighted educators than with BLV educators. Further analysis is found in Section~\ref{sec:november}. Nature of Context values are not in bold because it is a categorical variable.}
\end{table*}



% ChartGemma (3B) vs Qwen2 VL (2B) SFT on our dataset -> SightationCompletions 결과로 옮길것. 시각화 다시 할 것.
% TODO: Re-index the Experiment IDs
\begin{table*}\label{tab:big3c}
    \centering
    \begin{tabular}{llll}
    \toprule
    Experiment ID & & \multicolumn{2}{c}{\textbf{Assessments for}} \\ \cmidrule{3-4}
    \textbf{Description Generators}         & \textbf{Metrics} & $\textbf{Desc}^{\textcolor{ColorBlue}{\texttt{chartgemma}}}_{}$ & $\textbf{Desc}^{\textcolor{ColorOrange}{\texttt{q2bsft}}}_{}$  \\ \midrule

    \multirow{10}{0.2\textwidth}{Experiment 3c \\\textsc{\textcolor{ColorBlue}{ChartGemma (3B)}}\\vs.\\\textsc{\textcolor{ColorOrange}{Fine-tuned Qwen2-VL-2B-Instruct}}}
    
        & CLIP Score                                & 0.450 & \textbf{0.550} \\ 
        & SigLIP Score                              & 0.872 & \textbf{0.940} \\ 
        & BLIP-2 Retrieval Score                    & \textbf{0.511} & 0.490 \\ 
        & Self-BLEU                                 & \textbf{0.305} & 0.280 \\ 
        & PAC-Score                                 & 0.705 & \textbf{0.716}
        \\
        & LongClip-B                                & 0.316 & \textbf{0.684}
        \\
        & LongClip-L                                & \textbf{0.559} & 0.441\\ \cmidrule{2-4}                                       

        & $\cdot$ VLM-as-a-Judge Evaluation Average & 2.951 & \textbf{3.860} \\
        & Factuality                                & 3.068 & \textbf{4.119} \\
        & Informativeness                           & 2.848 & \textbf{3.967} \\
        & Succinctness                              & 3.253 & \textbf{3.925} \\
        & Diversity                                 & 2.635 & \textbf{3.428} \\

        % TODO: LVLM Judge Evaluation
        
    \bottomrule
    \end{tabular}
\caption{A 2B model fine-tuned on \textsc{SightationCompletions} outperforms a 3B model tuned on a larger dataset. Note that \textsc{ChartGemma} is not meant for conversational use. Hence, for a fair comparison, we did \textit{not} enter our guided generation prompt and instead input only the brief request ``Generate a caption'' to both models.}
\end{table*}

% BLIP2 baseline, BLIP2 coco-tuned, and BLIP2 Sightation-tuned comparison on a retrieval task
% with partial tuning and just 1 positive and 1 random negative,
% already improves many Recall and Precision values
% TODO: hard-negative training and observe performance
% TODO: visualize
\begin{table*}[t]
    \centering
    \begin{tabular}{l>{\columncolor{gray!15}}c>{\columncolor{gray!15}}ccc>{\columncolor{gray!15}}c>{\columncolor{gray!15}}c}
    % \begin{tabular}{lcccccc}
    \toprule
    & \multicolumn{6}{c}{2-way Cross-validation of \textbf{BLIP-2}} \\ \cmidrule{2-7}
    \textbf{Train set} & \multicolumn{2}{c}{N/A (Pre-trained)} & \multicolumn{2}{c}{COCO} & \multicolumn{2}{c}{\textsc{SightationRetrieval} (Ours)} \\ \cmidrule(lr){2-3}\cmidrule(lr){4-5}\cmidrule(lr){6-7}
    \textbf{Test set} & COCO & Ours & COCO & Ours & COCO & Ours \\ \midrule
    Recall@1                                & 0.171 & 0.048 & 0.185 & 0.033 & 0.180 & 0.076 \\ 
    Recall@5                                & 0.767 & 0.210 & 0.831 & 0.134 & 0.766 & 0.348 \\ 
    Recall@10                               & ---   & 0.340 & ---   & 0.229 & ---   & 0.549 \\ 
    Precision@1                             & 0.856 & 0.371 & 0.924 & 0.250 & 0.900 & 0.585 \\ 
    Precision@5                             & 0.767 & 0.324 & 0.831 & 0.204 & 0.766 & 0.535 \\
    Precision@10                            & ---   & 0.263 & ---   & 0.175 & ---   & 0.425 \\ 
    \bottomrule
    \end{tabular}
    \caption{\textsc{SightationRetrieval} shows promising potential as a challenging and effective training material for image-to-text retrievers. Two important observations can be made: the model trained on our set generalizes to COCO better than the other direction; our model performs on par with the model that was both trained and tested on COCO.  $K=10$ values are missing for tests with COCO, since its samples contain only 5 positives each.}
    \label{tab:retrieval}
\end{table*}


\begin{figure}[p!]
    \centering
    \includegraphics[width=0.95\linewidth]{figures/spider_retrieval_blip2coco_vs_blip2ours.png}
    \caption{Retrieval performance was measured with 2-way cross validation. On our test set (Left), the COCO-tuned \textsc{BLIP-2} generalizes poorly, whereas on the COCO test set (Right), the \textsc{SightationRetrieval}-tuned \textsc{BLIP-2} performs on par with the COCO-tuned \textsc{BLIP-2}.} 
    \label{fig:retrieval_tuning}
\end{figure}



\section{Annotator Demographics and Interviews}\label{app:demo}

\subsection{Demographics}

\subsubsection{BLV Educators}
Please refer to Table~\ref{tab:blv_educators_demo}.

\begin{table}[h!]
    \centering
    \tiny
    \begin{tabular}{lccccll}
    \toprule
    \textbf{ID} & \textbf{Sex} & \textbf{Age} & \makecell[c]{\textbf{Teaching}\\\textbf{Experience}\\\textbf{(years)}} & \makecell[c]{\textbf{Onset}\\\textbf{Age}} & \makecell[c]{\textbf{AI}\\\textbf{Use,}\\\textbf{Generic}} & \makecell[c]{\textbf{AI}\\\textbf{Use,}\\\textbf{Accessibility}} \\
    \midrule
    B1 & M & 54 & 28  & 16 & ChatGPT, Gemini & SenseReader \\
    B2 & F & 46 & 21  & Congenital & ChatGPT & SenseReader \\
    B3 & M & 47 & 5  & 9 & ChatGPT, Gemini & SenseReader \\
    B4 & M & 51 & 26 & 14 & SeeingAI, ChatGPT, Adot, Perplexity, Adot & SenseReader, NVDA, VoiceOver \\
    B5 & M & 20 &  1 & Congenital & SeeingAI, ChatGPT & SenseReader, NVDA \\
    B6 & M & 46 & 19 & --- & --- & SenseReader \\
    B7 & M & 44 & 21 &  Congenital & Be\_My\_Eyes, SeeingAI, ChatGPT, Claude & SenseReader, VoiceOver \\
    B8 & M & 45 & 19 &  Congenital & Be\_My\_Eyes, SeeingAI, ChatGPT & SenseReader, VoiceOver \\
    \bottomrule
    \end{tabular}
    \caption{BLV Teachers Information. All the BLV teachers in our study were of blindness level 1, the severest.}
    \label{tab:blv_educators_demo}
\end{table}

\subsubsection{Sighted Educators}
Please refer to Table~\ref{tab:sighted_educators_demo}.

\begin{table}[h!]
    \centering
    \footnotesize
    \begin{tabular}{lcccl}
    \toprule
    \textbf{ID} & \textbf{Sex} & \textbf{Age} & \textbf{Teaching Experience (years)} & \textbf{AI Use - Generic} \\
    \midrule
    S1 & M & 39 & 6.5 & ChatGPT \\
    S2 & M & 51 & 20 & ChatGPT, wrtn \\
    S3 & M & 48 & 21 & ChatGPT \\
    S4 & F & 40 & 13 & ChatGPT \\
    S5 & F & 56 & 33 & --- \\
    S6 & F & 49 & 20 & ChatGPT \\
    S7 & M & 49 & 20 & Gemini \\
    S8 & F & 49 & 24 & ChatGPT, Claude \\
    S9 & M & 44 & 14 & --- \\
    S10 & F & 50 & 20 & ChatGPT \\
    \bottomrule
    \end{tabular}
    \caption{Sighted Teachers Information.}
    \label{tab:sighted_educators_demo}
\end{table}




\section{Prompts}\label{app:prompts}

\centering
\begin{tcolorbox}[
    colback=white,
    colframe=gray!20,
    arc=10pt,
    boxrule=1pt,
    title={\textcolor{white}{\textbf{Prompts for Context Generation}}},
    coltitle=white,
    fonttitle=\large,
    colbacktitle=black!70,
]

\paragraph{\textbf{Generating} $\textbf{Desc}_{}$}
You are a helpful expert who is knowledgeable in various fields of academia. You are skilled in reading, interpreting, and understanding academic papers and figures contained therein. You are tasked with elaborating on the given information, which consists of a figure image. Write an informative and explanatory text in one paragraph under 200 words that describes the basic characteristics of the figure and incorporates important information. You may attempt to internally identify implicit points of curiosity for someone who is trying to understand the given figure, and then include explanations for those points in your response. Avoid mere reiteration of the given information as much as possible. You need not specify the origins of various parts of your response. [Optional: Aspect Suffix]

\paragraph{\textbf{Generating} $\textbf{Desc}_{\texttt{++}}$}
You are a helpful expert who is knowledgeable in various fields of academia. You are skilled in reading, interpreting, and understanding academic papers and figures contained therein. You are tasked with elaborating on the given information, which consists of a figure image and several question-answer pairs that have been derived from the figure. Write an informative and explanatory text in one paragraph under 200 words that describes the basic characteristics of the figure and incorporates important information from the question-answer pairs. You may attempt to internally identify implicit points of curiosity for someone who is trying to understand the given figure, and then include explanations for those points in your response. Avoid mere reiteration of the given information as much as possible. You need not specify the origins of various parts of your response. Here is the reference information: [QA\_PAIRS: {vqas}] [Optional: Aspect Suffix]

\end{tcolorbox}


\centering
\begin{tcolorbox}[
    colback=white,
    colframe=gray!20,
    arc=10pt,
    boxrule=1pt,
    title={\textcolor{white}{\textbf{Aspect Suffixes}}},
    coltitle=white,
    fonttitle=\large,
    colbacktitle=black!70,
]

\paragraph{Factuality}
When generating the diagram description, pay close attention to making it factual. A highly factual text delivers only the facts that are grounded in the diagram.

\paragraph{Informativeness}
When generating the diagram description, pay close attention to making it informative. A highly informative text describes all of the diagram, holistically.

\paragraph{Succinctness}
When generating the diagram description, pay close attention to making it succinct. A highly succinct text is concise and to the point.

\paragraph{Diversity}
When generating the diagram description, pay close attention to making it diverse. A highly diverse text captures a variety of perspectives from the diagram and employs multiple effective ways of getting the diagram message across. 

\end{tcolorbox}


% TODO: include augmented prompts per aspect.

\centering
\begin{tcolorbox}[
    colback=white,
    colframe=gray!20,
    arc=10pt,
    boxrule=1pt,
    title={\textcolor{white}{\textbf{Prompt for Question-answer Pair Generation}}},
    coltitle=white,
    fonttitle=\large,
    colbacktitle=black!70,
]

Please generate six question-and-answer pairs based on the provided image to aid in creating a comprehensive context. This context should include all essential details, allowing BLV (Blind and Low Vision) users to rely on the generated text instead of viewing the image (\textit{e.g.}, accessing information audibly). The question-and-answer pairs should cover both the main structure and finer details present in the image.
\end{tcolorbox}

\centering
\begin{tcolorbox}[
    colback=white,
    colframe=gray!20,
    arc=10pt,
    boxrule=1pt,
    title={\textcolor{white}{\textbf{Prompts for Reasoning Path Generation}}},
    coltitle=white,
    fonttitle=\large,
    colbacktitle=black!70,
]

You will be provided with a diagram, along with two descriptions of it. As an expert and experienced educator, you are tasked to examine your descriptions to identify common reasoning paths, such as cause-effect relationships, step-by-step processes, explanations of phenomena, comparisons of contrasts, and dependencies between components.

The identified reasoning paths should be under 25 words. Please provide the reasoning paths that you examined in the following JSON format:

\begin{verbatim}
        {"Context1": {"ReasoningPath": text},
         "Context2": {"ReasoningPath": text}}
\end{verbatim}

DO NOT return anything other than the JSON above.
\end{tcolorbox}

%\clearpage
\centering
\begin{tcolorbox}[
    colback=white,
    colframe=gray!20,
    arc=10pt,
    boxrule=1pt,
    title={\textcolor{white}{\textbf{Prompts for Reasoning QA Generation}}},
    coltitle=white,
    fonttitle=\large,
    colbacktitle=black!70,
]

You will be provided with a diagram, along with two descriptions of it. As an expert and experienced educator, you are tasked to examine your descriptions to generate reasoning question and answer pairs of five categories such as:

\begin{itemize}
    \item Causal Reasoning: "Why does [event] happen?"
    \item Process Reasoning: "What happens after [event]?"
    \item Conditional Reasoning: "What if [condition] changes?"
    \item Explanatory Reasoning: "Explain the role of [component] in the process."
    \item Reverse Reasoning: "Given [outcome], what might have caused it?"
\end{itemize}

Please provide the reasoning question and answer pairs that you generated in the following JSON format:

\begin{verbatim}
        {"Context 1": {"Causal": {"Question": text,
                                    "Answer": text},
                        "Process": {"Question": text,
                                    "Answer": text},	
                        "Conditional": {"Question": text,
                                        "Answer": text},
                        "Explanatory": {"Question": text,
                                        "Answer": text},
                        "Reverse": {"Question": text,
                                    "Answer": text}
            },
            "Context 2": {"Causal": {"Question": text,
                                    "Answer": text},
                        ..
                        "Reverse": {"Question": text,
                                    "Answer": text}
        }
\end{verbatim}

Each generated question should be under 15 words and each corresponding answer should be under 25 words.

DO NOT return anything other than the JSON above.

\end{tcolorbox}


% TODO: ChartGemma LLM judge prompt 참고했다 언급할 것. 
\centering
\begin{tcolorbox}[
    colback=white,
    colframe=gray!20,
    arc=10pt,
    boxrule=1pt,
    title={\textcolor{white}{\textbf{Prompt for VLM-as-a-Judge Evaluation of Description Pairs}}},
    coltitle=white,
    fonttitle=\large,
    colbacktitle=black!70,
]
You will be provided with a diagram, along with two descriptions of it. As an expert and experienced educator, you are tasked to evaluate each description based on the following qualities on a 5-point Likert scale. For each statement, give a score corresponding to how strongly you agree with the given statement: 1 (Strongly Disagree), 2 (Disagree), 3 (Neutral), 4 (Agree), or 5 (Strongly Agree).

\begin{itemize}
    \item Diversity: The description captures a variety of perspectives from the diagram and conveys multiple effective ways of getting the diagram message across.
    \item Succinctness: The description is concise and to the point, avoiding unnecessary details.
    \item Factuality: The description is accurate and reflects solely the information presented in the diagram.
    \item Informativeness: The description covers the diagram holistically, and it effectively conveys the main trends and insights of the diagram.
\end{itemize}

Please provide your ratings in the following JSON format:

\begin{verbatim}
        {
            Context 1: {
                Diversity: score,
                Succinctness: score,
                Factuality: score,
                Informativeness: score
            },
            Context 2: {
                Diversity: score,
                Succinctness: score,
                Factuality: score,
                Informativeness: score
            }
        }
\end{verbatim}

DO NOT return anything other than the JSON above.

\end{tcolorbox}



\begin{tcolorbox}[
    colback=white,
    colframe=gray!20,
    arc=10pt,
    boxrule=1pt,
    title={\textcolor{white}{\textbf{Prompt for VLM-as-a-Judge Evaluation of Question-Answer Pairs}}},
    coltitle=white,
    fonttitle=\large,
    colbacktitle=black!70,
]
\textbf{Instruction}

You need to rate the quality of the given Question and Answer in relation to a diagram. Specifically, assess whether the Q\&A correctly references and interprets the information presented in the diagram. Consider the Q\&A’s clarity, specificity, and coherence as they pertain to the diagram’s content.You must take into account not only the Answer but also whether an appropriate Question has been provided for the given diagram at the same time.
The rating scale is as follow:

\begin{itemize}
    \item very poor: The Q\&A is unclear, vague, or incoherent in relation to the diagram. It lacks essential information or misinterprets the diagram’s content.
    \item poor: The Q\&A is somewhat unclear or omits important details from the diagram. It requires significant clarification or correction to align with the diagram.
    \item average: The Q\&A is moderately clear and specific. It may need additional details or minor clarifications to fully match the diagram’s information.
    \item good: The Q\&A is clear, specific, and mostly well-formed in referencing the diagram. It provides sufficient context to understand how the diagram supports the question and answer.
    \item excellent: The Q\&A is very clear, specific, and well-articulated. It precisely references and fully aligns with the diagram, containing all necessary details and context.
\end{itemize}

\textbf{Output Format}


Given the user’s diagram, question, and answer, you must:

Provide an assessment that briefly explains the strengths and/or weaknesses of how the Q\&A relates to the diagram.
Output your rating (one of: very poor, poor, average, good, excellent) by filling in the placeholders below.

\begin{verbatim}
[
{
	"explanation": "[...]",
	"input_quality": "[very poor/poor/average/good/excellent]"
},
...
]
\end{verbatim}

\textbf{Notes}

\begin{itemize}
\item DO NOT return anything else other than the JSON above.
\item Number of item in above list should be same as the number of given QA pairs. Also the order for the explanation and input quality should be same as input QA's order
\end{itemize}

\end{tcolorbox}

\section{Fine-tuning Configurations}\label{app:tuning_config}


\begin{table*}[t!]
\centering
\small
\begin{tabularx}{\textwidth}{lXX}
\toprule
\textbf{Parameter} & \textbf{SFT Config (Qwen2-VL-2B-Instruct)} & \textbf{DPO Config (Qwen2-VL-2B-Instruct)} \\ 
\midrule
\multicolumn{3}{l}{\textbf{Script Arguments}} \\
Dataset Name & \textsc{SightationCompletions} & \textsc{SightationPreference} \\
\midrule
\multicolumn{3}{l}{\textbf{Training Configurations}} \\
Output Directory & anonymous & anonymous \\
Evaluation Strategy & steps & steps \\
Train Batch Size & 1 & 1 \\
Evaluation Batch Size & 1 & 1 \\
Gradient Accumulation Steps & 8 & 8 \\
Training Epochs & 1 & 1 \\
Save Total Limit & 5 & 5 \\
bfloat16 Enabled & true & true \\
Evaluation Steps & 10 & 10 \\
Label Names & ["labels"] & ["labels"] \\
Load Best Model at End & true & true \\
Metric for Best Model & eval\_loss & eval\_loss \\
Use Liger & true & true \\
Max Sequence Length & 1024 & 1024 \\
Remove Unused Columns & false & true \\
Dataset Kwargs & skip\_prepare\_dataset: true & skip\_prepare\_dataset: false \\
Gradient Checkpointing & true & true \\
Gradient Checkpointing Kwargs & use\_reentrant: false & use\_reentrant: false \\
Dataset Num Processors & 8 & 8 \\
Torch Compile & true & --- \\
DDP Find Unused Parameters & --- & true \\
\midrule
\multicolumn{3}{l}{\textbf{Model Config}} \\
Use PEFT & false & false \\
Model Path & Qwen/Qwen2-VL-2B-Instruct & Qwen/Qwen2-VL-2B-Instruct \\
Torch Dtype & bfloat16 & bfloat16 \\
Attention Implementation & flash\_attention\_2 & flash\_attention\_2 \\
\bottomrule
\end{tabularx}
\caption{SFT and DPO configurations for Qwen2-VL-2B-Instruct. Tuning was performed on 4 \texttimes A6000 GPUs.
}
\label{tab:tuning_config_2b_sft_dpo}
\end{table*}

\begin{table*}[t]
\centering
\small
\begin{tabularx}{\textwidth}{lXX}
\toprule
\textbf{Parameter} & \textbf{SFT Config (Qwen2-VL-7B-Instruct)} & \textbf{DPO Config (Qwen2-VL-7B-Instruct)} \\ 
\midrule
\multicolumn{3}{l}{\textbf{Script Arguments}} \\
Dataset Name & \textsc{SightationCompletions}  &  \textsc{SightationPreference} \\
\midrule
\multicolumn{3}{l}{\textbf{Training Configurations}} \\
Output Directory & anonymous & anonymous \\
Evaluation Strategy & steps & steps \\
Train Batch Size & 1 & 1 \\
Evaluation Batch Size & 1 & 1 \\
Gradient Accumulation Steps & 8 & 8 \\
Training Epochs & 1 & 1 \\
Save Total Limit & 5 & 5 \\
bfloat16 Enabled & true & true \\
Evaluation Steps & 10 & 10 \\
Label Names & ["labels"] & ["labels"] \\
Load Best Model at End & false & false \\
Metric for Best Model & eval\_loss & eval\_loss \\
Use Liger & true & true \\
Max Sequence Length & 1024 & 1024 \\
Remove Unused Columns & false & true \\
Dataset Kwargs & skip\_prepare\_dataset: true & skip\_prepare\_dataset: false \\
Gradient Checkpointing & true & true \\
Gradient Checkpointing Kwargs & use\_reentrant: false & use\_reentrant: false \\
Dataset Num Processors & 8 & 8 \\
DDP Find Unused Parameters & true & true \\
\midrule
\multicolumn{3}{l}{\textbf{Model Config}} \\
Use PEFT & true & true \\
Model Path & Qwen/Qwen2-VL-7B-Instruct & Qwen/Qwen2-VL-7B-Instruct \\
Torch Dtype & bfloat16 & bfloat16 \\
Attention Implementation & flash\_attention\_2 & flash\_attention\_2 \\
LoRA Rank (r) & 16 & 16 \\
LoRA Alpha & 16 & 16 \\
LoRA Dropout & 0.1 & 0.1 \\
LoRA Target Modules & q\_proj, k\_proj, v\_proj, o\_proj, gate\_proj, up\_proj, down\_proj & q\_proj, k\_proj, v\_proj, o\_proj, gate\_proj, up\_proj, down\_proj \\
\bottomrule
\end{tabularx}
\caption{SFT and DPO configurations for Qwen2-VL-7B-Instruct. Tuning was performed on 4 \texttimes A6000 GPUs.
}
\label{tab:tuning_config_7b_sft_dpo}
\end{table*}


\begin{table*}[t]
\centering
\small
\begin{tabularx}{\textwidth}{lX}
\toprule
\textbf{Component} & \textbf{Configuration} \\ \midrule
\textbf{Model} & BLIP-2 (Salesforce/blip2-itm-vit-g) \\
\textbf{GPUs} & Text model on CUDA:0, Vision model on CUDA:1 \\
\textbf{Dataset} & \textsc{SightationRetrieval} \\
\textbf{Loss} & InfoNCE (temperature = 0.07) \\
\textbf{Batch Size} & 1 (with gradient accumulation steps = 4) \\
\textbf{Epochs} & 5 \\
\textbf{Optimizer} & AdamW (Text LR: 5e-5, Vision LR: 2e-5) \\
\textbf{Gradient Clipping} & Max norm = 1.0 \\
\textbf{Scheduler} & Linear warmup (10\% of steps) \\
\textbf{Frozen Layers} & All except: layernorm, projection, encoder layers 10-11 (Vision); layernorm, projection, encoder layers 10-11, crossattention (Text) \\
\textbf{Checkpoints} & Best and per-epoch saved to anonymized path \\
\bottomrule
\end{tabularx}
\caption{Training configurations for BLIP-2 image-text retrieval.}
\label{tab:retrieval_config}
\end{table*}



% TODO: include fine tuning details. see frozen YAMLs in Run Configurations.

\clearpage
\section{Guidelines}\label{app:guidelines}


\centering
\begin{tcolorbox}[
    colback=white,
    colframe=gray!20,
    arc=10pt,
    boxrule=1pt,
    title={\textcolor{white}{\textbf{Annotation Guidelines for the Sighted General Group (1/4)}}},
    coltitle=white,
    fonttitle=\large,
    colbacktitle=black!70,
]

Please carefully read the guidelines below and ensure accurate labeling. Your responses are considered high-quality data and can have critical implications for the experiment. Pay special attention to the Caution section.
% TODO: structure for better readability -> Done.

\paragraph{Annotation Guidelines}
Thank you for contributing to this project. In the following paragraphs, we will walk you through the project description, your tasks, and annotation examples.

\paragraph{Project}
Our project targets the visually impaired. People who are Blind or have Low Vision (BLV) do not always benefit from the latest AI developments in the same way or extent as sighted users. In this pilot study, we would like to first assess exactly how much state-of-the-art models may assist sighted users, so that we may gain insights into (i) what state-of-the-art models can and cannot do and (ii) what modifications might be necessary to alter their assistive information to cater specifically for BLV users.

\paragraph{Task}
Each task you are about to complete consists of:
\begin{itemize}
    \item 1 image
    \item 2 image description pairs, each containing two texts
\end{itemize}
Given these, you are tasked with:
\begin{itemize}
    \item Selecting the overall “winner” for each pair
    \item Rating the qualities of each text
    \item Copying and pasting the best contributing sentence for each text
\end{itemize}

\centering
\includegraphics[width=0.5\linewidth, keepaspectratio]{figures/label_studio_sample.png} 
           
\end{tcolorbox}


\centering
\begin{tcolorbox}[
    colback=white,
    colframe=gray!20,
    arc=10pt,
    boxrule=1pt,
    title={\textcolor{white}{\textbf{Annotation Guidelines for the Sighted General Group (2/4)}}},
    coltitle=white,
    fonttitle=\large,
    colbacktitle=black!70,
]

Here is a detailed instruction on each task.

\paragraph{Selecting the overall winner}
\begin{itemize}
    \item Have a look at both texts. In your opinion, which is the better description of the text?
    \item This may be based on general impression or your internal criteria.
\end{itemize}

    \centering
    \includegraphics[width=0.95\linewidth]{figures/label_studio_sample2.png}

\end{tcolorbox}


\centering
\begin{tcolorbox}[
    colback=white,
    colframe=gray!20,
    arc=10pt,
    boxrule=1pt,
    title={\textcolor{white}{\textbf{Annotation Guidelines for the Sighted General Group (3/4)}}},
    coltitle=white,
    fonttitle=\large,
    colbacktitle=black!70,
]

\paragraph{Rating the qualities of each text}
\begin{itemize}
    \item We break down what makes an image description a “good” one into a few qualities.
    \item Do not feel overly pressured to justify or rationalize your choice of the overall winner. In fact, you may treat the “overall winner” choice task completely independent of the quality rating task.
    \item On a scale from 1 to 5, how well does the text exhibit the following qualities?
    \begin{itemize}
        \item Factuality:
        \begin{itemize}
            \item Does the text contain facts about the image content?
            \item You may give a low score if the text contains statements that cannot be inferred from the image. These extraneous claims can include knowledge from the “world external to the image”, even though that knowledge itself may be true.
            \item You may give a low score if the text contains wrong statements about the image.
        \end{itemize}
        \item Informativeness:
        \begin{itemize}
            \item Does the text holistically describe the image content and help you become better informed about it?
            \item You may give a low score if some parts of the image seem “left out” in the text description.
        \end{itemize}
        \item Succinctness:
        \begin{itemize}
            \item Does the text describe the image content in a concise yet helpful way?
            \item Judgments based solely on text length should be avoided.
            \item Instead, think of the “density” of information contained in the text.
            \item You may give a low score if the text contains redundant/repeated information, inefficient sentence structures, and/or overly simple vocabulary that tend to make the text feel “sparse”.
        \end{itemize}
        \item Diversity:
        \begin{itemize}
            \item Does the text help you understand the image in various ways?
            \item There may be multiple effective descriptors about one image. There may be multiple perspectives and/or approaches to understanding one image. Do you think the given text addresses these?
            \item You may give a low score if the text feels too focused on small parts of the image or views the image in an overly specific (possibly contrived) perspective to lay out the description.
        \end{itemize}
    \end{itemize}
\end{itemize}

\begin{center}
    \includegraphics[width=0.95\linewidth]{figures/label_studio_sample3.png}    
\end{center}

\end{tcolorbox}



\centering
\begin{tcolorbox}[
    colback=white,
    colframe=gray!20,
    arc=10pt,
    boxrule=1pt,
    title={\textcolor{white}{\textbf{Annotation Guidelines for the Sighted General Group (4/4)}}},
    coltitle=white,
    fonttitle=\large,
    colbacktitle=black!70,
]


\paragraph{Copying and pasting the best sentence for each text}
\begin{itemize}
    \item For each text, drag (highlight) the sentence that has best contributed to each text overall.
    \item The best sentence does not have to function as a one-sentence summary.
    \item This should be from the English text, not the Korean translation.
    \item Paste the sentence into the text field and press enter to submit.
    \item You will see your submitted best sentence on the display. Press the trash can button (Delete) to change your mind and submit a different sentence.
\end{itemize}

\begin{center}
    \includegraphics[width=0.95\linewidth]{figures/label_studio_sample4.png}    
\end{center}

\paragraph{Caution}
\begin{itemize}
    \item Text length should not be a criterion for your assessment.
    \item Best sentence should be copied from the English paragraph.
\end{itemize}

\end{tcolorbox}



\centering
\begin{tcolorbox}[
    colback=white,
    colframe=gray!20,
    arc=10pt,
    boxrule=1pt,
    title={\textcolor{white}{\textbf{Evaluation Guidelines for the \textcolor{lime}{Sighted Educator} Group (1/2)}}},
    coltitle=white,
    fonttitle=\large,
    colbacktitle=black!70,
]

Thank you again for joining our experiments.

Attached is a spreadsheet containing 40 images, each of which has two descriptions, produced by various AI models in response to our request to generate a context for the input diagram.

As annotators, you are tasked with evaluating the quality of these texts.

\paragraph{Selecting the Preferred Text}

Have a look at the diagram in the ``Image'' column, along with the two contexts written in the ``Context1'' and ``Context2'' columns. In the ``Preferred Text'' column, enter your choice as 1 or 2.
\begin{itemize}
    \item This preference may rely on your own personal criteria.
\end{itemize}

\end{tcolorbox}




\centering
\begin{tcolorbox}[
    colback=white,
    colframe=gray!20,
    arc=10pt,
    boxrule=1pt,
    title={\textcolor{white}{\textbf{Evaluation Guidelines for the \textcolor{lime}{Sighted Educator} Group (2/2)}}},
    coltitle=white,
    fonttitle=\large,
    colbacktitle=black!70,
]


\paragraph{Quantitative Assessment}

Give a score ranging from 1 to 5 for each quality listed below. For each statement, enter your assessment on a scale from a 5 if you ``Strongly Agree'' to a 1 if you ``Strongly Disagree''.

\begin{itemize}
    \item Factuality: The text delivers only facts that are grounded on the diagram.
    \begin{itemize}
        \item Even if a piece of knowledge in the text is factual, you may give a low score if that knowledge cannot be inferred from the diagram.
        \item Wrong textual descriptions of the diagram content also have low merit.
    \end{itemize}
    \item Informativeness: The text describes all of the diagram, holistically.
    \begin{itemize}
        \item You may give a low score to a context leaving out parts of the diagram.
    \end{itemize}
    \item Succinctness: The text is concise and to the point.
    \begin{itemize}
        \item Please assess whether the context conveys an appropriate ``density'' of information.
        \item You may give a low score if a context seems repetitive.
        \item Please avoid scoring based on apparent text length.
    \end{itemize}
    \item Diversity: The text captures a variety of perspectives from the diagram and employs multiple effective ways of getting the diagram message across.
    \begin{itemize}
        \item There may be multiple different ways to understand a diagram. Please assess whether these ways have been put together in the given context.
    \end{itemize}
    \item Usefulness: The text is helpful to BLV.
    \begin{itemize}
        \item As an experienced educator for learners with visual impairments, please evaluate how useful and helpful the text would be.
    \end{itemize}
\end{itemize}

\paragraph{Qualitative Assessment}

In the ``Reason'' column, please justify your preference choice (\textit{i.e.}, the 1 or 2 selection) with a brief explanation. Simple comments, as long as they tell us the textual quality your choice was based on, may still prove helpful for our research. See examples below.

\begin{itemize}
    \item ``contains various descriptions of ants''
    \item ``written more logically''
    \item ``Context 1 contains a more realistic definition of a food web.''
    \item ``2 is more concise.''
    \item ``Context 1 lacks a description of the artery.''
    \item ``While both texts faithfully address the rotation and revolution movements of the Earth, Context 1 describes in-depth how they manifest as different natural phenomena on the planet.''
\end{itemize}

\end{tcolorbox}



\centering
\begin{tcolorbox}[
    colback=white,
    colframe=gray!20,
    arc=10pt,
    boxrule=1pt,
    title={\textcolor{white}{\textbf{Evaluation Guidelines for the \textcolor{Yellow}{BLV Educator} Group}}},
    coltitle=white,
    fonttitle=\large,
    colbacktitle=black!70,
]

Thank you again for joining our experiments.

Attached is a spreadsheet containing 80 images, each with a description that has been produced by various AI models in response to our request to generate a context for the input diagram.\\

The spreadsheet comes with 81 rows and 8 columns. The table headers are as follows: Image, Context, Succinctness, Diversity, Usefulness (Summary), Usefulness (Multiple-choice Questions), Usefulness (Open-ended Questions), and Nature of Context. Apart from the header row, the remaining 80 rows each contain 1 image and 1 description.\\

As annotators, you are tasked with evaluating these texts, based on the text alone.

\paragraph{Quantitative Assessment}

Give a score ranging from 1 to 5 for each quality listed below. For each statement, on a scale from a 5 if you ``Strongly Agree'' to a 1 if you ``Strongly Disagree''.

\begin{itemize}
    \item Succinctness: The text is concise and to the point.
    \begin{itemize}
        \item Please assess whether the context conveys an appropriate ``density'' of information.
        \item You may give a low score if a context seems repetitive.
        \item Please avoid scoring based on apparent text length.
    \end{itemize}
    \item Diversity: The text captures a variety of perspectives from the diagram and employs multiple effective ways of getting the diagram message across.
    \begin{itemize}
        \item There may be multiple different ways to understand a diagram. Please assess whether these ways have been put together in the given context.
    \end{itemize}
    \item Usefulness (Summary): The text serves as a good summary.
    \begin{itemize}
        \item Please assess how well the text helps you formulate an idea of the diagram content.
    \end{itemize}
    \item Usefulness (Multiple-choice Questions): The text would be useful in solving short-answer, multiple-choice questions based on the diagram.
    \begin{itemize}
        \item Suppose you are to solve short-answer multiple-choice questions that have been constructed from the diagram. How well would the context help you answer these questions?
    \end{itemize}
    \item Usefulness (Open-ended Questions): The text would be useful in solving descriptive, essay type questions based on the diagram.
    \begin{itemize}
        \item Suppose you are to answer an open-ended long-answer question that has been constructed from the diagram. How well would the context help you answer such a question?
    \end{itemize}
    \item Nature of Context: The text is rich in interpretive detail.
    \begin{itemize}
        \item On a scale of 1 to 5, if the text appears to lay out plain and straightforward facts from the diagram, give a score of 1. If it rather contains interpretive descriptions, and/or reasoned explanations, give a score of 5.
    \end{itemize}
\end{itemize}

\end{tcolorbox}
